% original_pdf: source/普通高考/2017/2017全国3理(云南,广西,贵州,四川,西藏).pdf, page 1, question 7.
% evidence: pdfimages -list found no embedded bitmap; the source flowchart is vector line art.
% The original page was rendered at 300 dpi; 100px coarse and 50px fine grids locate the
% comparison crop in tmp/redraw_flowchart_2017_national_paper_3_science/.
% Type: program flowchart (schematic; topology and labels are the mathematical content).
% elements: rounded start/end terminals, input/output trapezoids, three assignment rectangles,
%           the t<=N decision diamond, vertical arrows, the 否 output branch, and the left loop.
% relations: 开始 -> 输入 N -> t=1,M=100,S=0 -> t<=N; 是 -> S=S+M -> M=-M/10 -> t=t+1,
%            then return to the decision; 否 -> 输出 S -> 结束.
% solid segments: every node outline, vertical connector, branch connector, and arrowhead.
% dashed segments: none.
% labels: formulas are centered in process nodes; 是 is below the diamond, 否 is above/right
% of the output branch; the return path stays left of the process column and re-enters at the top.
% The initial stem and return path meet before the decision; only one downward arrow enters t<=N,
% while the loop's horizontal segment retains its right-pointing arrowhead.

\begin{tikzpicture}[
    x=0.9cm,
    y=1.2cm,
    >=Stealth,
    line width=0.45pt,
    line cap=butt,
    line join=miter,
    flowtext/.style={font=\normalsize,anchor=center,align=center,
                     execute at begin node={\setlength{\parindent}{0pt}\noindent}},
    every node/.style={flowtext},
    flowbox/.style={draw,flowtext,inner xsep=4pt,inner ysep=2.5pt},
    terminal/.style={flowbox,rounded corners=0.18cm,minimum width=1.40cm,minimum height=0.68cm},
    process/.style={flowbox,minimum height=0.68cm},
    io/.style={flowbox,trapezium,trapezium left angle=72,trapezium right angle=108,
               minimum height=0.74cm},
    decision/.style={flowbox,diamond,aspect=2.75,inner xsep=4pt,inner ysep=2.5pt,
                     minimum height=0.96cm},
]
    % Keep a narrow, even safety margin around the outer loop and output branch.
    \path[use as bounding box] (-4.80,0.35) rectangle (5.85,10.40);

    % The clean schematic coordinates preserve the source's relative layout.
    \node[terminal] (start) at (0,9.85) {开始};
    \node[io,minimum width=2.20cm] (input) at (0,8.52) {输入\(N\)};
    \node[process,minimum width=4.75cm] (init) at (0,7.18)
        {\(t=1,M=100,S=0\)};
    \node[decision,minimum width=4.00cm] (test) at (0,5.70) {\(t\leq N\)};
    \node[process,minimum width=2.90cm] (sum) at (0,4.18) {\(S=S+M\)};
    \node[process,minimum width=2.80cm,minimum height=1.45cm] (scale) at (0,2.54)
        {\(M=-\dfrac{M}{10}\)};
    \node[process,minimum width=2.50cm] (increment) at (0,0.91) {\(t=t+1\)};
    \node[io,minimum width=2.70cm] (output) at (4.02,4.18) {输出\(S\)};
    \node[terminal] (finish) at (4.02,2.84) {结束};

    % Main vertical path and the affirmative branch.
    \draw[->] (start.south) -- (input.north);
    \draw[->] (input.south) -- (init.north);
    % The initial stem joins the loop junction without an extra arrowhead;
    % one downward arrow is shared by both paths into the decision node.
    \draw (init.south) -- (0,6.55);
    \draw[->] (test.south) -- node[right=2pt] {是} (sum.north);
    \draw[->] (sum.south) -- (scale.north);
    \draw[->] (scale.south) -- (increment.north);

    % The negative branch bypasses the loop body and reaches 输出 S.
    \coordinate (output-turn) at (4.02,5.70);
    \draw[->] (test.east) -- node[flowtext,above=2pt] {否} (output-turn)
        -- (output.north);
    \draw[->] (output.south) -- (finish.north);

    % Return path from t=t+1 to the top of the decision diamond.
    \coordinate (loop-left) at (-4.28,0.91);
    \coordinate (loop-top) at (-4.28,6.55);
    \draw (increment.west) -- (loop-left) -- (loop-top);
    % The source shows a right-pointing arrow terminating at the loop junction
    % and a separate downward arrow into the decision node.
    \draw[->] (loop-top) -- (0,6.55);
    \draw[->] (0,6.55) -- (test.north);
\end{tikzpicture}
